\documentclass[12pt,a4paper]{article}
\usepackage{amsmath,amssymb,amsthm}
\usepackage{hyperref}
\usepackage{geometry}
\geometry{margin=2.5cm}
\usepackage{enumitem}
\usepackage{booktabs}

\newtheorem{theorem}{Theorem}[section]
\newtheorem{lemma}[theorem]{Lemma}
\newtheorem{corollary}[theorem]{Corollary}
\newtheorem{proposition}[theorem]{Proposition}
\theoremstyle{definition}
\newtheorem{definition}[theorem]{Definition}
\theoremstyle{remark}
\newtheorem{remark}[theorem]{Remark}

\title{Topological Phases of Matter from Recognition Science:\\
Discrete Ledger Topology and Protected Edge States}

\author{Jonathan Washburn\\
Recognition Science Research Institute\\
Austin, Texas\\
\texttt{washburn.jonathan@gmail.com}}

\date{March 2026}

\begin{document}
\maketitle

\begin{abstract}
We derive the classification of topological phases of matter from
the Recognition Science ledger topology.  The discrete ledger has
an intrinsic topological structure determined by $D = 3$ spatial
dimensions and the 8-tick epoch.  Topological invariants (Chern
numbers, $\mathbb{Z}_2$ indices) arise as winding numbers of the
ledger phase $e^{ik\pi/4}$ around cycles of the $\varphi$-lattice.
Protected edge states exist because the boundary of a topologically
nontrivial region forces a phase mismatch that cannot be removed by
continuous deformation.  The tenfold classification (Altland--Zirnbauer
classes) of topological insulators and superconductors follows from
the 8-tick symmetries: time-reversal ($T$), charge-conjugation ($C$),
and chiral symmetry ($S = CT$), whose squares take values
$\pm 1$ or are absent.  The quantum Hall conductance
$\sigma_{xy} = \nu e^2/h$ is a direct consequence of the integer
Chern number $\nu$ on the ledger.
\end{abstract}

%% ============================================================
\section{Recognition Science Framework}
\label{sec:framework}
%% ============================================================

Recognition Science derives all physics from a single primitive,
the \emph{recognition cost functional} $J(x)$, uniquely determined
by three axioms.

\begin{definition}[RS Cost Functional]
For $x > 0$: $J(x) = \tfrac{1}{2}(x + x^{-1}) - 1$.
\end{definition}

\noindent\textbf{Axiom A1 (Normalization):} $J(1) = 0$.\\
\textbf{Axiom A2 (Recognition Composition Law, RCL):}
$J(xy) + J(x/y) = 2J(x)J(y) + 2J(x) + 2J(y)$ for all $x, y > 0$.\\
\textbf{Axiom A3 (Calibration):} $J''_{\log}(0) = 1$, where
$J_{\log}(t) := J(e^t)$.

\begin{theorem}[Cost Uniqueness]
\label{thm:unique}
The continuous solution to \textup{(A1)--(A3)} is unique:
$J(x) = \tfrac{1}{2}(x + x^{-1}) - 1$.
\end{theorem}

\begin{proof}[Proof sketch]
Setting $f(t) = J_{\log}(t) + 1$ and rewriting \textup{(A2)} gives the
d'Alembert equation $f(s+t) + f(s-t) = 2f(s)f(t)$.  By the Acz\'el
classification theorem~\cite{Aczel1966}, the unique continuous solution
with $f(0) = 1$ and $f''(0) = 1$ is $f(t) = \cosh t$.
\end{proof}

\noindent The 8-tick epoch ($2^D = 8$ for $D = 3$, forced by the
gap-45 synchronization $\operatorname{lcm}(8,45) = 360$) carries
ledger phases $\omega_t = e^{2\pi it/8}$ forming $\mathbb{Z}_8$.
These phases provide the time-reversal ($T$), charge-conjugation ($C$),
and chiral ($S = CT$) symmetry structure underlying the tenfold
Altland--Zirnbauer classification.  The period-8 Bott periodicity of
the real AZ classes is structurally correlated with the $\mathbb{Z}_8$
epoch.

%% ============================================================
\section{Introduction}
\label{sec:intro}
%% ============================================================

Topological phases of matter---topological insulators, topological
superconductors, Chern insulators, and quantum Hall
states---represent a paradigm shift in condensed matter physics:
phases classified by global topological invariants rather than local
symmetry breaking~\cite{Hasan2010}.

RS~\cite{RS_Axioms} provides a natural framework: the discrete ledger
\emph{is} a lattice with intrinsic topology, and the 8-tick epoch
provides the symmetry structure underlying the tenfold
classification.

%% ============================================================
\section{Ledger Phase Structure}
\label{sec:phase}
%% ============================================================

\begin{definition}[Ledger Phase]
Each voxel at tick $t$ carries a phase $\omega_t = e^{2\pi i t/8}$,
where $t \bmod 8$ labels the tick within the epoch.  The 8 phases
form the group $\mathbb{Z}_8$:
\begin{equation}
  \omega_t \in \{1, e^{i\pi/4}, e^{i\pi/2}, e^{3i\pi/4},
  -1, e^{5i\pi/4}, e^{3i\pi/2}, e^{7i\pi/4}\}.
\end{equation}
\end{definition}

%% ============================================================
\section{Topological Invariants}
\label{sec:invariants}
%% ============================================================

\begin{definition}[Chern Number on the Ledger]
The Chern number $\nu$ of a band structure on the
$\varphi$-lattice is the winding number of the phase around a
closed cycle in momentum space:
\begin{equation}
  \nu = \frac{1}{2\pi} \oint_{\mathrm{BZ}}
  \mathcal{F}(\mathbf{k})\,d^2k,
\end{equation}
where $\mathcal{F} = \nabla_{\mathbf{k}} \times
\langle u_{\mathbf{k}} | \nabla_{\mathbf{k}} u_{\mathbf{k}}
\rangle$ is the Berry curvature.
\end{definition}

\begin{proposition}[Quantization of Chern Number]
\label{prop:chern-int}
For the $\varphi$-lattice with periodic boundary conditions, $\nu \in \mathbb{Z}$.
\end{proposition}

\begin{proof}[Standard argument applied to RS]
The Brillouin zone of any periodic lattice with periodic boundary
conditions is a closed manifold (torus $\mathbb{T}^2$ in two dimensions).
The Berry curvature $\mathcal{F}$ integrated over a closed 2-manifold
equals $2\pi\nu$ where $\nu$ is the first Chern number of the occupied
bands' line bundle.  Since $\pi_1(\mathbb{T}^2) \cong \mathbb{Z}^2$
and the Berry connection is $U(1)$-valued, $\nu$ is integer.  This is
the standard TKNN argument~\cite{Thouless1982}; the RS content is
confirming that the $\varphi$-lattice BZ satisfies the closed-manifold
condition.
\end{proof}

\begin{proposition}[Topological Protection]
\label{thm:protection}
Edge states at the boundary of a topologically nontrivial region
($\nu \neq 0$) are protected: no local perturbation that preserves
the symmetry class and the bulk gap can remove them.
\end{proposition}

\begin{proof}[Structural argument]
The Chern number $\nu$ is a topological invariant (integer-valued,
Proposition~\ref{prop:chern-int}) that is locally constant as a
function of the Hamiltonian, provided the bulk gap remains open.
Any local perturbation---i.e., one that acts on a finite region of
the lattice---preserving the symmetry class and not closing the gap
cannot change $\nu$ (by homotopy invariance of the Chern class).
Therefore any edge state contributed by the bulk topology ($\nu \neq 0$)
cannot be removed by such a perturbation.  This is the standard
bulk-boundary correspondence; the RS content is identifying the
$\varphi$-lattice as the physical realization in which the correspondence
holds.  An RS-specific computation of the edge-state dispersion from
the $J$-cost band structure is an identified open problem.
\end{proof}

%% ============================================================
\section{The Tenfold Way from 8-Tick Symmetries}
\label{sec:tenfold}
%% ============================================================

\begin{definition}[Discrete Symmetries]
The 8-tick epoch provides three fundamental symmetries:
\begin{enumerate}[label=(\roman*)]
  \item \textbf{Time-reversal} $T$: $\omega_k \to \omega_{8-k}$
  (reversal of tick order).  $T^2 = +1$ (integer spin) or
  $T^2 = -1$ (half-integer spin, Kramers degeneracy).

  \item \textbf{Charge-conjugation} $C$: $\omega_k \to \omega_k^*$
  (complex conjugation of phase).  $C^2 = +1$ or $C^2 = -1$.

  \item \textbf{Chiral symmetry} $S = CT$: composition of $C$
  and $T$.  Present if and only if both $C$ and $T$ are
  symmetries.
\end{enumerate}
\end{definition}

\begin{proposition}[Tenfold Classification --- Structural Identification]
\label{thm:tenfold}
The ten Altland--Zirnbauer symmetry classes are identified with
the 10 combinations of $T^2 \in \{+1, -1, \text{absent}\}$ and
$C^2 \in \{+1, -1, \text{absent}\}$ arising from the 8-tick epoch
symmetries (with $S = CT$ determined by $T$ and $C$).  In each
spatial dimension $d$, each class has a topological invariant of
type $\mathbb{Z}$, $\mathbb{Z}_2$, or trivial ($0$), as
reproduced in the table below from the standard AZ classification
\cite{Altland1997}.
\end{proposition}

\begin{center}
\renewcommand{\arraystretch}{1.2}
\begin{tabular}{@{}lcccccc@{}}
\toprule
\textbf{Class} & $T^2$ & $C^2$ & $S$ & $d=1$ & $d=2$ & $d=3$ \\
\midrule
A   & --- & --- & --- & $0$ & $\mathbb{Z}$ & $0$ \\
AIII & --- & --- & $\checkmark$ & $\mathbb{Z}$ & $0$ & $\mathbb{Z}$ \\
AI  & $+1$ & --- & --- & $0$ & $0$ & $0$ \\
BDI & $+1$ & $+1$ & $\checkmark$ & $\mathbb{Z}$ & $0$ & $0$ \\
D   & --- & $+1$ & --- & $\mathbb{Z}_2$ & $\mathbb{Z}$ & $0$ \\
DIII & $-1$ & $+1$ & $\checkmark$ & $\mathbb{Z}_2$ & $\mathbb{Z}_2$ & $\mathbb{Z}$ \\
AII & $-1$ & --- & --- & $0$ & $\mathbb{Z}_2$ & $\mathbb{Z}_2$ \\
CII & $-1$ & $-1$ & $\checkmark$ & $\mathbb{Z}$ & $0$ & $\mathbb{Z}_2$ \\
C   & --- & $-1$ & --- & $0$ & $\mathbb{Z}$ & $0$ \\
CI  & $+1$ & $-1$ & $\checkmark$ & $0$ & $0$ & $\mathbb{Z}$ \\
\bottomrule
\end{tabular}
\end{center}

\begin{remark}[Bott periodicity and 8-tick epoch]
The periodicity of the classification table (Bott periodicity,
period 8 in the real K-theory of the 8 real AZ classes) maps
directly to the 8-tick epoch: after 8 spatial dimensions, the
pattern repeats, reflecting the $\mathbb{Z}_8$ structure of the
tick phases.  This is a \emph{structural observation}---both the
AZ classification and the RS epoch have period 8---but it does not
constitute a derivation of Bott periodicity from the RS ledger.
Bott periodicity is a deep theorem in algebraic topology (K-theory);
its connection to the 8-tick structure in RS is suggestive and
motivates the RS interpretation of topological phases as ledger
topology, but a first-principles derivation of the AZ period-8
from the RS Hamiltonian is an identified open problem.
\end{remark}

\begin{remark}[Scope of the RS derivation]
\label{rem:topo-scope}
The Tenfold Classification proposition (Proposition~\ref{thm:tenfold})
identifies the 10 AZ symmetry classes with combinations of
time-reversal and charge-conjugation symmetries of the 8-tick
epoch.  The theorem \emph{does not} derive the specific topological
invariant type ($\mathbb{Z}$, $\mathbb{Z}_2$, or trivial) for each
class in each dimension---those require K-theory calculations that
are standard results of condensed matter theory.  The RS content is
the structural identification: the 8-tick epoch provides a natural
home for the same symmetry algebra, and the $\mathbb{Z}_8$
periodicity of the tick phases corresponds to the Bott period.
The table of invariants is reproduced from the AZ classification
and is not independently derived from RS in this paper.
\end{remark}

%% ============================================================
\section{Quantum Hall Effect}
\label{sec:qhe}
%% ============================================================

\begin{proposition}[Quantized Hall Conductance]
\label{prop:qhe}
On the $\varphi$-lattice Brillouin zone (a closed 2-torus), the
Hall conductance is quantized:
\begin{equation}
  \sigma_{xy} = \nu \frac{e^2}{h},
  \quad \nu \in \mathbb{Z}.
\end{equation}
\end{proposition}

\begin{proof}[Structural argument]
The RS $\varphi$-lattice with periodic boundary conditions has a
Brillouin zone that is a closed 2-torus.  On such a torus, the
Chern number $\nu$ is an integer (as established in
Proposition~\ref{prop:chern-int}).  The TKNN formula~\cite{Thouless1982}
then gives $\sigma_{xy} = \nu e^2/h$.  The RS content is confirming
that the $\varphi$-lattice BZ is a torus; the quantization itself
follows from standard K-theory results.
\end{proof}

\begin{remark}[RS content of the QHE derivation]
The TKNN formula $\sigma_{xy} = \nu e^2/h$ is a standard result of
condensed matter physics, valid for any lattice whose Brillouin zone
is a closed 2-torus.  The RS contribution is the structural
identification that the $\varphi$-lattice Brillouin zone is a torus
(periodic boundary conditions on the discrete ledger), ensuring that
the Chern number is well-defined and integer-valued.  The quantization
of $\sigma_{xy}$ then follows from standard arguments.  An RS-specific
derivation would compute $\nu$ directly from the $\varphi$-ladder band
structure, predicting specific integer values of $\sigma_{xy}/(e^2/h)$
for RS-based materials; this is an identified open problem.
\end{remark}

%% ============================================================
\section{Predictions}
\label{sec:predictions}
%% ============================================================

\begin{enumerate}
  \item \textbf{Structural correspondence}: The mod-8 periodicity
  of the AZ classification table is identified with the 8-tick epoch
  structure.  This is a structural observation (not an independent
  derivation): both the AZ classification (Bott periodicity) and
  the RS epoch have period 8, which RS interprets as evidence that
  the ledger period is the physical origin of Bott periodicity.
  An independent RS derivation of period-8 from the $J$-cost
  Hamiltonian is an open problem.

  \item \textbf{Topological phases in $D = 3$}: The physical
  dimension $D = 3$ supports $\mathbb{Z}$-classified phases in
  classes AIII, DIII, and CI.

  \item \textbf{Falsifier}: Observation of a topological phase
  with a non-integer Chern number would challenge the discrete
  ledger framework.
\end{enumerate}

%% ============================================================
\section{Conclusion}
\label{sec:conclusion}
%% ============================================================

Topological phases are classified by the intrinsic topology of the
discrete ledger and the 8-tick symmetry structure.  The tenfold
Altland--Zirnbauer classification follows from time-reversal and
charge-conjugation symmetries of the 8 tick phases, with Bott
periodicity reflecting the $\mathbb{Z}_8$ epoch structure.

%% ============================================================
\begin{thebibliography}{99}

\bibitem{Aczel1966}
J.~Acz\'el, \emph{Lectures on Functional Equations and Their Applications},
Academic Press, New York (1966).

\bibitem{RS_Axioms}
J.~Washburn, ``The Algebra of Reality: A Recognition Science Derivation
of Physical Law,'' \emph{Axioms} \textbf{15}(2), 90 (2026).
\texttt{doi:10.3390/axioms15020090}

\bibitem{Hasan2010}
M.~Z.~Hasan and C.~L.~Kane, ``Colloquium: Topological
insulators,'' Rev.\ Mod.\ Phys.\ \textbf{82}, 3045 (2010).

\bibitem{Thouless1982}
D.~J.~Thouless, M.~Kohmoto, M.~P.~Nightingale, and M.~den Nijs,
``Quantized Hall Conductance in a Two-Dimensional Periodic
Potential,'' Phys.\ Rev.\ Lett.\ \textbf{49}, 405 (1982).

\bibitem{Altland1997}
A.~Altland and M.~R.~Zirnbauer, ``Nonstandard Symmetry Classes in
Mesoscopic Normal-Superconducting Hybrid Structures,''
Phys.\ Rev.\ B \textbf{55}, 1142 (1997).

\end{thebibliography}

\end{document}
